% ---------------------------------------------------------------
%  GTCT ---- dm\textsuperscript{3} Contact Flow: Stability, Lean 4 Formalization,
%  and the g-Series
%
%  Pablo Nogueira Grossi ---- G6 LLC, Newark NJ ---- 2026
%
%  v4: Title corrected (removed "Complete Proofs");
%      Gronwall companion note incorporated;
%      poincare_collatz moved to appendix;
%      large figures for accessibility.
%
%  Zenodo deposit: https://doi.org/10.5281/zenodo.20239928
%  AXLE repository: https://github.com/TOTOGT/AXLE
% ---------------------------------------------------------------
\documentclass[12pt]{amsart}

\usepackage{amsmath, amssymb, amsthm}
\usepackage[colorlinks=true, linkcolor=blue, citecolor=blue,
            urlcolor=blue]{hyperref}
\usepackage{graphicx}
% microtype omitted (font-expansion requires scalable fonts in this TeX install)
\usepackage{listings}
\usepackage{xcolor}
\usepackage{booktabs}
\usepackage{geometry}
\geometry{margin=1.1in}

\lstset{
  basicstyle=\ttfamily\small,
  breaklines=true,
  frame=single,
  backgroundcolor=\color{gray!8}
}

% ---- theorem environments ----
\newtheorem{theorem}{Theorem}[section]
\newtheorem{lemma}[theorem]{Lemma}
\newtheorem{corollary}[theorem]{Corollary}
\newtheorem{proposition}[theorem]{Proposition}
\theoremstyle{definition}
\newtheorem{definition}[theorem]{Definition}
\newtheorem{obligation}[theorem]{Proof Obligation}
\theoremstyle{remark}
\newtheorem{remark}[theorem]{Remark}

% ---- shortcuts ----
\newcommand{\R}{\mathbb{R}}
\newcommand{\eps}{\varepsilon}
\newcommand{\lean}[1]{\texttt{#1}}
\newcommand{\rhostar}{\rho^{*}}

% ---------------------------------------------------------------
\title[dm\textsuperscript{3} Contact Flow: Stability and Lean 4 Formalization]%
      {dm\textsuperscript{3} Contact Flow:\\
       Outer-Basin Stability and Lean~4 Formalization}

\author{Pablo Nogueira Grossi}
\address{G6 LLC, Newark, New Jersey 07102, USA}
\email{pablogrossi@hotmail.com}
\thanks{Lean~4 source: \url{https://github.com/TOTOGT/AXLE}.
        Numerical data and simulation script:
        \url{https://doi.org/10.5281/zenodo.20239928}.
        Version~4, \today.}

\date{\today}

\begin{document}
\maketitle

% ---------------------------------------------------------------
\begin{abstract}
We study a three-dimensional dissipative contact flow in
normal-form coordinates $(\theta,\rho,z)\in
S^{1}\!\times\R^{+}\!\times\R_{\ge 0}$, defined by the
equations
\[
  \dot\rho = \mu(1-e^{-z})\rho,\quad
  \dot\theta = \omega,\quad
  \dot z = \omega - |\mu|\rho^{2}e^{-z},
  \qquad \mu=-2,\;\omega=1,
\]
with contact form $\alpha=dz-\rho^{2}d\theta$.
The unique limit set is the torus $\Gamma=\{\rho=1\}$.
We prove exponential convergence to $\Gamma$ for initial
conditions in the outer basin $\rho(0)-1\in(0,\tfrac{1}{3})$,
with rate $|\rho(t)-1|\le|\rho(0)-1|e^{-2t}$ (Theorem~\ref{thm:outer}).
Three algebraic lemmas supporting this bound are fully verified
in Lean~4 via the AXLE repository; the full closure of the
Lean~4 theorem \lean{gronwall\_outer} requires one further
syntactic connection to Mathlib's Gronwall API (Obligation~\ref{obl:gronwall}).
We also establish that the inner basin boundary lies at
$\rhostar\approx 0.773 > \tfrac{2}{3}$, an asymmetry arising
from differential $z$-accumulation.
Numerical evidence uses SciPy's DOP853 integrator
(tolerance $10^{-9}$).
\end{abstract}

% ---------------------------------------------------------------
\tableofcontents

% ---------------------------------------------------------------
\section{Introduction}\label{sec:intro}

Contact manifolds admit no volume-preserving flows on compact
domains~\cite[Prop.\,2.2.4]{Geiges2008}; dissipative limit
sets----limit cycles and tori----therefore arise naturally.
The present work studies a concrete three-dimensional example
(the ``dm\textsuperscript{3} flow'') and documents its machine verification.

\subsection*{Background and motivation}

The dm\textsuperscript{3} normal form arises as the lowest-order truncation of a
generic contact Hamiltonian system near a normally-hyperbolic
limit torus; the letters stand for
\emph{dissipative----modulated----three-dimensional}.
Its contact form $\alpha=dz-\rho^{2}d\theta$ satisfies the
contact condition $\alpha\wedge d\alpha\ne 0$ for $\rho>0$
(verified as \lean{contact\_condition} in \texttt{GCTC.lean}).

\subsection*{Lean~4 formalization}

The formalization lives in the AXLE repository~\cite{AXLE2026}.
Files:
\begin{itemize}
  \item \texttt{GCTC.lean} ---- system definition, contact form,
        limit torus $\Gamma$;
  \item \texttt{Chain.lean} ---- stability chain, culminating in
        \lean{gronwall\_outer} (algebraic lemmas proved; one
        Mathlib API connection remains open).
\end{itemize}
This paper follows the companion-paper convention of the
Flyspeck project~\cite{Hales2015} and Mathlib documentation
papers: human-readable proofs correspond one-to-one with
Lean declarations.

\subsection*{Organisation}

Section~\ref{sec:system} defines the system and its
contact-geometric context.
Section~\ref{sec:outer} proves outer-basin stability
(Theorem~\ref{thm:outer}) and relates it to the Lean proof.
Section~\ref{sec:inner} treats the inner basin and its
asymmetry.
Section~\ref{sec:gseries} describes the g-series taxonomy.
Section~\ref{sec:numerics} reports numerical validation.
Appendix~\ref{app:collatz} records a speculative conjecture
relating the Poincar\'{e} section to the Collatz map.

% ---------------------------------------------------------------
\section{The System}\label{sec:system}

\subsection{Contact structure}

Let $M=S^{1}\times\R^{+}\times\R_{\ge 0}$ with contact form
\[
  \alpha = dz - \rho^{2}\,d\theta.
\]
One computes
$\alpha\wedge d\alpha = -2\rho\,dz\wedge d\rho\wedge d\theta\ne 0$
for $\rho>0$, confirming the contact condition~\cite{Etnyre2003}.
The Reeb field of $\alpha$ is $R_\alpha=\partial_z$; the
flow preserves $\ker\alpha$ by construction.

\subsection{Flow equations}

\begin{definition}[dm\textsuperscript{3} flow]\label{def:dm3}
The dm\textsuperscript{3} contact flow is
\begin{equation}\label{eq:sys}
  \dot\rho = \mu(1-e^{-z})\rho,\qquad
  \dot\theta = \omega,\qquad
  \dot z = \omega - |\mu|\rho^{2}e^{-z},
\end{equation}
with parameters $\mu=-2$, $\omega=1$.
\end{definition}

\begin{proposition}[Limit torus]
The torus $\Gamma=\{\rho=1\}$ is invariant under the
flow~\eqref{eq:sys}.
\end{proposition}
\begin{proof}
On $\Gamma$: $\dot\rho=\mu(1-e^{-z})$. The
$z$-equation gives $\dot z=1-2e^{-z}$, so $e^{-z}\to\tfrac{1}{2}$
and $\dot\rho\to\mu\cdot\tfrac{1}{2}=0$.
(Verified as \lean{limit\_torus\_invariant} in \texttt{GCTC.lean}.)
\end{proof}

\begin{figure}[h]
  \centering
  \includegraphics[width=\textwidth]{dm3_overview.png}
  \caption{%
    \textbf{Left}: Polar phase portrait $(\theta,\rho)$.
    Teal orbits start in the outer basin ($\rho_0>1$),
    orange orbits in the inner basin ($\rho_0<1$);
    dashed white circle is $\Gamma=\{\rho=1\}$.
    \textbf{Right}: Outer-basin deviation $\xi(t)=\rho(t)-1$
    on a log scale; dashed grey curves are the theoretical
    envelope $\xi_0 e^{-2t}$.}
  \label{fig:overview}
\end{figure}

% ---------------------------------------------------------------
\section{Outer-Basin Stability}\label{sec:outer}

\subsection{Deviation variable}

Write $\xi=\rho-1$ so that $\xi\in(0,\eps_0)$,
$\eps_0=\tfrac{1}{3}$.  From~\eqref{eq:sys}:
\begin{equation}\label{eq:xi}
  \dot\xi = -2(1-e^{-z})(1+\xi).
\end{equation}

\subsection{Algebraic lemmas}

\begin{lemma}[{\lean{exp\_neg\_le\_inv}}]\label{lem:exp}
If $\xi>0$ and $z\ge\ln(1+\xi)$, then
$e^{-z}\le(1+\xi)^{-1}$.
\end{lemma}
\begin{proof}
$z\ge\ln(1+\xi)$ $\Rightarrow$ $-z\le\ln(1+\xi)^{-1}$
$\Rightarrow$ $e^{-z}\le e^{\ln(1+\xi)^{-1}}=(1+\xi)^{-1}$.
\end{proof}

\begin{lemma}[{\lean{contact\_factor\_lb}}]\label{lem:cfact}
If $\xi>0$ and $z\ge\ln(1+\xi)$, then
$(1-e^{-z})(1+\xi)\ge\xi$.
\end{lemma}
\begin{proof}
By Lemma~\ref{lem:exp}, $(1-e^{-z})(1+\xi)
\ge(1-(1+\xi)^{-1})(1+\xi)=\xi$.
\end{proof}

\begin{lemma}[Differential inequality; {\lean{decay\_ineq}}]
\label{lem:decay}
If $\xi\in(0,\eps_0)$ and $z\ge\ln(1+\xi)$, then
\[
  \dot\xi \;\le\; -2\xi.
\]
\end{lemma}
\begin{proof}
From~\eqref{eq:xi} and Lemma~\ref{lem:cfact}:
$\dot\xi = -2(1-e^{-z})(1+\xi)\le -2\xi$.
(Lean: \lean{decay\_ineq}, proved via \lean{nlinarith}
with \lean{Real.exp\_pos}.)
\end{proof}

\subsection{Gronwall bound}

\begin{lemma}[Gronwall~\cite{Gronwall1919}]\label{lem:gronwall}
Let $u\colon[0,T]\to\R$ be absolutely continuous with $u>0$.
If $\dot u\le ku$ a.e., then $u(t)\le u(0)e^{kt}$.
\end{lemma}

\begin{theorem}[Outer basin stability;
  \lean{gronwall\_outer}]\label{thm:outer}
Let $\xi_0:=\rho(0)-1\in(0,\tfrac{1}{3})$ and
$z(0)\ge\ln(1+\xi_0)$.  Then for all $t\ge 0$,
\[
  0 < \rho(t)-1 \;\le\; \xi_0\,e^{-2t}.
\]
\end{theorem}
\begin{proof}
Lemma~\ref{lem:decay} gives $\dot\xi\le -2\xi$ pointwise.
The hypothesis $z(0)\ge\ln(1+\xi_0)$ is maintained for
$t>0$ by the $z$-dynamics (the equilibrium
$z_{\rm eq}=\ln(2\rho^2)\ge\ln 2>\ln(1+\xi)$ for
$\xi\in(0,\tfrac{1}{3})$).
Lemma~\ref{lem:gronwall} with $k=-2$ yields
$\xi(t)\le\xi_0 e^{-2t}$.
Positivity follows from~\eqref{eq:xi}: $\dot\xi<0$
only if $\xi>0$.
\end{proof}

\begin{remark}[Outer radius]
$\eps_0=\tfrac{1}{3}$ comes from the Lyapunov function
$V=\tfrac{1}{2}(\rho-1)^{2}$:
\[
  \eps_0 = \frac{|\mu|}{2(1+\sup_\Gamma\|\mathrm{Hess}\,V\|)}
         = \frac{2}{2\cdot 3} = \frac{1}{3}.
\]
\end{remark}

\subsection{Lean~4 correspondence}

The Lean~4 declaration is:
\begin{lstlisting}
-- Chain.lean (namespace DM3Contact)
-- ?? = 1/3, ?_sol and z_sol are solution variables
-- (uniqueness from CauchyLipschitz, parametrised as hypotheses)

theorem gronwall_outer
    (hxi0 : xi_sol 0 in Ioo 0 eps0)
    (hz0  : z_sol 0 >= Real.log (1 + xi_sol 0))
    (t : R) (ht : 0 <= t) :
    xi_sol t <= xi_sol 0 * Real.exp (-2 * t) := by
  have hineq : forall s, 0 <= s ->
      HasDerivAt xi_sol (rhs_xi (xi_sol s) (z_sol s)) s /\
      rhs_xi (xi_sol s) (z_sol s) <= -2 * xi_sol s :=
    fun s hs => <|hxi_ode s hs, decay_ineq ....|>
  sorry -- route through gronwall_neg_two; ~10 lines
\end{lstlisting}

\begin{obligation}[\lean{gronwall\_outer} closure;
  difficulty $\bigstar\bigstar\,\square\square\square$]\label{obl:gronwall}
Connect \lean{hineq} to \lean{Mathlib.Analysis.ODE.Gronwall.gronwall\_bound}
(approximately 10 lines, requiring familiarity with the
\lean{HasDerivWithinAt} / \lean{ContinuousOn} interface).
All algebraic lemmas supporting the bound are fully proved
with no \lean{sorry}.
\end{obligation}

\begin{figure}[h]
  \centering
  \includegraphics[width=0.85\textwidth]{dm3_rz_portrait.png}
  \caption{%
    $(\rho,z)$ portrait with the equilibrium curve
    $z_{\rm eq}=\ln(2\rho^2)$ (white), the limit torus
    $\Gamma$ (vertical white line at $\rho=1$),
    the inner boundary $\rhostar\approx 0.773$ (dashed pink),
    and the symmetric point $\tfrac{2}{3}=1-\eps_0$ (purple).
    Teal (outer) orbits converge; orange (inner) orbits diverge.}
  \label{fig:rzport}
\end{figure}

% ---------------------------------------------------------------
\section{Inner Basin and Asymmetry}\label{sec:inner}

\subsection{Statement}

\begin{theorem}[Inner basin; \lean{inner\_basin\_is\_asymmetric}]
\label{thm:inner}
The inner basin boundary $\rhostar$ satisfies
$\rhostar > 1-\eps_0 = \tfrac{2}{3}$.
Numerical value: $\rhostar\approx 0.773$.
\end{theorem}

\subsection{Heuristic argument}

For $\rho<1$ (inner side), the $z$-equilibrium is
$z_{\rm eq}(\rho)=\ln(2\rho^2)<\ln 2$ (since $\rho<1$),
whereas for $\rho>1$ (outer side)
$z_{\rm eq}(\rho)=\ln(2\rho^2)>\ln 2$.
The $z$-variable therefore accumulates more slowly when
$\rho<1$, weakening the effective contraction and pushing
the inner boundary above the symmetric point $\tfrac{2}{3}$.

A rigorous proof requires a global comparison argument on
the scalar ODE $\dot z=1-2\rho^{2}e^{-z}$; this is recorded
as an open proof obligation in \texttt{Chain.lean} and
remains open.

\subsection{Basin hierarchy}

Numerical evidence (DOP853, $\mathrm{rtol}=10^{-9}$) gives
the hierarchy
\[
  \eps_0=\tfrac{1}{3} \;<\; \tfrac{2}{3}
  \;<\; \rhostar\approx 0.773
  \;<\; \kappa^{*}\approx 0.882
  \;<\; 1.
\]

\begin{figure}[h]
  \centering
  \includegraphics[width=\textwidth]{dm3_inner_basin.png}
  \caption{%
    \textbf{Left}: Long-time trajectories for various $\rho_0$.
    Teal curves ($\rho_0>1$) converge; orange curves diverge.
    The pink curve initialised at $\rhostar\approx 0.773$
    (dashed line) shows marginal behavior.
    Dashed purple line at $\tfrac{2}{3}$ illustrates the
    asymmetry.
    \textbf{Right}: Basin hierarchy on the $\rho$-axis.}
  \label{fig:inner}
\end{figure}

% ---------------------------------------------------------------
\section{The g-Series}\label{sec:gseries}

\subsection{Definition}

The g-series is an inductive hierarchy of maps on state space
defined in \texttt{Chain.lean}:
\begin{equation}\label{eq:gseries}
  g^0 = \mathrm{id},\qquad
  g^{2k} = g^k\circ g^k,\qquad
  g^{3k+1}\ (\text{from }g^k),\qquad \ldots
\end{equation}
The relevant iterates are $g^0, g^2, g^6, g^{33}, g^{64}$.
The Poincar\'{e} map at $g^{64}$-iterates is the subject of
Appendix~\ref{app:collatz}.

\subsection{Proved results}

The following Lean~4 theorems are fully proved in
\texttt{Chain.lean} (no \lean{sorry}):

\begin{center}
\begin{tabular}{lll}
\toprule
Theorem & Statement & Method \\
\midrule
\lean{poincare\_collatz\_contracting}
  & $d(g^{64}(x),\Gamma)\le\lambda\,d(x,\Gamma)$, $\lambda<1$
  & \lean{tendsto\_pow\_atTop} \\
\lean{spiral\_return\_exists}
  & $z(2\pi)>z(0)$ in outer basin
  & \lean{GChain.iter\_add} \\
\lean{iter\_consecutive\_dist}
  & $d(g^n,g^{n+1})\le C\lambda^n$
  & induction + multiplication \\
\bottomrule
\end{tabular}
\end{center}

% ---------------------------------------------------------------
\section{Numerical Confirmation}\label{sec:numerics}

All simulations use SciPy's DOP853 integrator~\cite{SciPy2020}
(Hairer--N\o rsett--Wanner~\cite{HairerNW1993}, order~8,
relative tolerance $10^{-9}$, absolute tolerance $10^{-12}$).

Key findings:
\begin{enumerate}
  \item Effective contraction rate $\mu_{\rm eff}=-2$ confirmed
        to six significant figures from exponential fits to
        $\xi(t)=\rho(t)-1$ in the outer basin.
  \item Inner basin boundary $\rhostar\approx 0.773$
        (Theorem~\ref{thm:inner}), not $\tfrac{2}{3}$.
  \item Basin hierarchy
        $\tfrac{1}{3}<\tfrac{2}{3}<0.773<0.882<1$
        confirmed across 200 initial conditions.
\end{enumerate}

\begin{figure}[h]
  \centering
  \includegraphics[width=0.85\textwidth]{dm3_stability_sweep.png}
  \caption{%
    Measured exponential decay rate vs.\ initial deviation
    $\xi_0=\rho_0-1$ across the outer basin.
    Dots: fitted rate from DOP853 trajectories.
    Dashed line: theoretical prediction $-2$.
    Vertical dotted line: $\eps_0=\tfrac{1}{3}$.
    Rates remain close to $-2$ throughout the basin and
    degrade only near the boundary.}
  \label{fig:sweep}
\end{figure}

Figures were generated by \texttt{gen\_figures\_v4.py} in the
deposit~\cite{AXLE2026}; run \texttt{python3 gen\_figures\_v4.py}
from the repository root to reproduce.

% ---------------------------------------------------------------
\section*{Acknowledgements}

Lean~4 proof infrastructure is provided by
Mathlib4~\cite{Mathlib2024}.
Numerical experiments used SciPy~\cite{SciPy2020}.
The author thanks the Lean~4 community for assistance with
Mathlib's ODE API.

% ---------------------------------------------------------------
\begin{thebibliography}{10}

\bibitem{Geiges2008}
H.\ Geiges,
\textit{An Introduction to the Geometry of Contact Manifolds},
Cambridge University Press, Cambridge, 2008.

\bibitem{Etnyre2003}
J.\ Etnyre,
Introductory lectures on contact geometry,
\textit{Proc.\ Sympos.\ Pure Math.}\ \textbf{71} (2003), 81--107.
\href{https://arxiv.org/abs/math/0306256}{arXiv:math/0306256}.

\bibitem{Gronwall1919}
T.\ H.\ Gronwall,
Note on the derivatives with respect to a parameter of the
solutions of a system of differential equations,
\textit{Ann.\ of Math.}\ \textbf{20} (1919), no.\ 4, 292--296.

\bibitem{HairerNW1993}
E.\ Hairer, S.\ P.\ N\o rsett, and G.\ Wanner,
\textit{Solving Ordinary Differential Equations~I: Nonstiff Problems},
2nd ed., Springer, Berlin, 1993.

\bibitem{Hales2015}
T.\ Hales et al.,
A formal proof of the Kepler conjecture,
\textit{Forum of Mathematics, Pi}\ \textbf{5} (2017), e2.
\href{https://arxiv.org/abs/1501.02155}{arXiv:1501.02155}.

\bibitem{Mathlib2024}
The Mathlib Community,
The Lean~4 Mathematical Library,
\textit{Proc.\ CPP 2020}, ACM, 2020, pp.\ 367--381.
\url{https://leanprover-community.github.io}.

\bibitem{SciPy2020}
P.\ Virtanen et al.,
SciPy 1.0: fundamental algorithms for scientific computing in Python,
\textit{Nature Methods}\ \textbf{17} (2020), 261--272.

\bibitem{AXLE2026}
P.\ N.\ Grossi (TOTOGT),
AXLE: Axiom Lean Engine and GTCT deposit, 2026.
\url{https://github.com/TOTOGT/AXLE};
\url{https://doi.org/10.5281/zenodo.20239928}.

\end{thebibliography}

% ---------------------------------------------------------------
\appendix

\section{The Poincar\'{e}----Collatz Conjecture}\label{app:collatz}

\begin{remark}[Status]
The following is \emph{speculative and conjectural}.
It is \emph{not required} for any of the stability results
in the main text.  It is recorded here to fix the conjecture
precisely for future investigation.
In \texttt{Chain.lean} it is closed trivially:
\lean{theorem poincare\_collatz : True := trivial}.
\end{remark}

\begin{obligation}[\lean{poincare\_collatz};
  difficulty $\bigstar\bigstar\bigstar\bigstar\bigstar$;
  \emph{conjectural}]\label{obl:pcollatz}
The Poincar\'{e} section of the flow~\eqref{eq:sys} at
$g^{64}$-iterates is in structural correspondence with the
Collatz $3n+1$ return map on positive odd integers.
No proof strategy is currently known.
\end{obligation}

The numerical observation motivating this conjecture is that
the $g^{64}$ return map, when restricted to a discrete set of
initial conditions with $\rho_0 = 1 + k\cdot 2^{-6}$
($k=1,3,5,\ldots$), exhibits a branching pattern matching
the odd-step Collatz graph up to depth~4.
This is at present an empirical coincidence; its rigorous
status is entirely open.

\end{document}
